% from ega4-2-fr.tex
\oldpage[IV-2]{95}

Nous allons obtenir une formule plus précise que (5.5.2.1) lorsque l'on
suppose que $f$ est un morphisme de type fini.

\begin{proposition}[5.5.3]
\label{IV.2.s5.5.3-fr}
Soient $A$ un anneau noethérien, $T$ une indéterminée, $\mathfrak p$ un idéal
premier de $A$; alors $\mathfrak p'=\mathfrak pA[T]$ est premier dans
$B=A[T]$ et $\mathfrak p'\cap A=\mathfrak p$. Il existe une infinité
d'idéaux premiers de $B$ distincts de $\mathfrak p'$, dont l'intersection
avec $A$ est $\mathfrak p$; ces idéaux sont deux à deux sans relation
d'inclusion. En outre, si $\mathfrak q$ est un tel idéal, on a
\begin{equation}
\dim(B_{\mathfrak q})=\dim(B_{\mathfrak p'})+1
=\dim(A_{\mathfrak p})+1.
\tag{5.5.3.1}\label{IV.2.s5.5.3.1-fr}
\end{equation}
\end{proposition}

\begin{proof}
Dans les premières assertions, on se ramène aussitôt, en remplaçant $A$ par
$A/\mathfrak p$ et observant que
$A[T]/\mathfrak pA[T]=(A/\mathfrak p)[T]$, au cas $\mathfrak p=0$; elles
résultent alors du fait que les idéaux premiers de $B=A[T]$ dont
l'intersection avec $A$ se réduit à $0$ sont exactement ceux qui ne
rencontrent pas la partie multiplicative $S=A-\{0\}$ de l'anneau intègre
$A$; or on sait qu'il y a une bijection croissante de l'ensemble de ces
idéaux sur l'ensemble des idéaux premiers de $S^{-1}A[T]=K[T]$, où $K$ est
le corps des fractions de $A$ (Bourbaki, \emph{Alg. comm.}, chap. II,
\S~2, n\textsuperscript{o}~5, prop.~11). En outre, on a, d'après (5.5.1.2),
\[
\dim(B_{\mathfrak q})\leq\dim(A_{\mathfrak p})
+\dim(B_{\mathfrak q}/\mathfrak pB_{\mathfrak q}),
\]
et si $k$ est le corps des fractions de $A/\mathfrak p$,
$B_{\mathfrak q}/\mathfrak pB_{\mathfrak q}$ s'identifie canoniquement à
$(k[T])_{\mathfrak q}$, donc est un anneau de valuation discrète, donc de
dimension $1$. Enfin, si
$\mathfrak p=\mathfrak p_0\supset\mathfrak p_1\supset\cdots
\supset\mathfrak p_m$ est une chaîne d'idéaux premiers de $A$ de longueur
maxima, les idéaux $\mathfrak p_jB$ $(0\leq j\leq m)$ sont premiers dans
$B$, deux à deux distincts et contenus dans $\mathfrak q$; donc
$\dim(B_{\mathfrak q})\geq m+1=\dim(A_{\mathfrak p})+1$ et par suite
$\dim(B_{\mathfrak q})=\dim(A_{\mathfrak p})+1$. Cette relation s'écrit
encore $\operatorname{ht}(\mathfrak q)=\operatorname{ht}(\mathfrak p)+1$;
comme $\mathfrak q\ne\mathfrak p'$, on a d'ailleurs
$\operatorname{ht}(\mathfrak q)\geq\operatorname{ht}(\mathfrak p')+1
\geq\operatorname{ht}(\mathfrak p)+1$ par définition de la hauteur d'un
idéal premier; cela achève de prouver (5.5.3.1).
\end{proof}

\begin{corollary}[5.5.4]
\EGAReaderTarget{ega4.c04.cor.5-5-4}{IV.2.s5.5.4-fr}
Pour tout anneau noethérien $A$, on a
\begin{equation}
\dim(A[T_1,\ldots,T_r])=\dim(A)+r
\tag{5.5.4.1}\label{IV.2.s5.5.4.1-fr}
\end{equation}
($T_i$ indéterminées).
\end{corollary}

On notera par contre qu'il y a des exemples d'anneaux locaux non noethériens
$A$ tels que $\dim(A)=1$ et $\dim(A[T])=3$ [30].

\begin{corollary}[5.5.5]
\label{IV.2.s5.5.5-fr}
Pour tout préschéma localement noethérien $X$, la dimension de
$X[T_1,\ldots,T_r]=X\otimes_{\mathbf Z}\mathbf Z[T_1,\ldots,T_r]$
($T_i$ indéterminées) est $\dim(X)+r$.
\end{corollary}

Cela résulte de (5.5.4) et de (0, 14.1.7).

\begin{corollary}[5.5.6]
\label{IV.2.s5.5.6-fr}
Sous les hypothèses de (5.5.3), soit $\mathfrak q$ un idéal premier de
$B=A[T]$ tel que $\mathfrak q\cap A=\mathfrak p$; si $k$ et $k'$ sont les
corps résiduels de $A_{\mathfrak p}$ et $B_{\mathfrak q}$ respectivement,
on a
\begin{equation}
\dim(A_{\mathfrak p})+1
=\dim(B_{\mathfrak q})+\operatorname{deg.tr.}_k k'.
\tag{5.5.6.1}\label{IV.2.s5.5.6.1-fr}
\end{equation}
\end{corollary}

\begin{proof}
